\section{Full spectrum after stationary-density regularization}\label{densityspectrum}
We now control the reconstructed chain on the regularized side of the
comparison left open in Section~\ref{regularization}. The result is uniform
over every admissible input density; no code algebra or lower bound for its
potential is imposed. It does not compare the reconstructed chain back to
the original chain.

Fix $0<\rho,b<1$, and put $s=(1+\rho)/(1-\rho)$ and $B=2^{1/s}$.
The constants in this section are fixed before the dimension, density,
matching order and smoothing parameter. Define
\begin{align*}
 q_0&=2\rho^{-2}+\rho^{-3}-1,& Q_*&=8s,\\
 \sigma&=\max\left\{4s,\frac{2Q_*^2}{(1-\rho^4)(q_0-1)}\right\},&
 F_\sigma(q)&=1+\frac{\sigma q/(\sigma+q)-1}{\rho^4}.
\end{align*}
Starting at $q_0$, set $q_{n+1}=F_\sigma(q_n)$ only until its first
crossing $q_{n_*}\geq Q_*$. Take $n_*=0$ if $q_0\geq Q_*$.
The proof below establishes that this crossing is finite. Set
\begin{align*}
 A_*&=b^{-1}(B/b)^{n_*},& C_*&=b^{-1}BA_*^2,\\
 \theta_{\rm in}&=\min\{1,\sqrt{s-1}\},&
 \theta_{\rm out}&=\min\{1,(s-1)^{-1/2}\},\\
 \tau_*&=\tfrac12\min\left\{\theta_{\rm in}\theta_{\rm out},
                    \sqrt{\frac{s-1}{\sigma-1}}\right\},&
 \gamma&=\frac1{2s-1}.
\end{align*}
In particular, $\tau_*>0$ and $\tau_*\leq1/2$. For an integer $k\geq0$,
write
\begin{equation}\label{densityspectrumerror}
 e_k=B\rho^{(k+1)/2},\qquad
 \mathcal E_k=(2e_k)^\gamma(C_*+B)^{1-\gamma}
                 +(B+1)2^{-(k+1)}+e_k^2.
\end{equation}

\begin{theorem}[Uniform full spectrum on the regularized side]\label{densityspectrumtheorem}
Let $L\geq1$, $0\leq k\leq L$, and let $\pi$ be a positive mean-one
density on the cube with
\[
 \pi\geq b,\qquad \E\pi^s\leq2,\qquad
 \E\pi\chi_S=0\quad(1\leq |S|\leq k).
\]
For every $0\leq\tau\leq\tau_*$, put $w=P_\tau\pi$.
There is a unique positive $u$ with $uPu=w$. The actual Doob kernel
$K_u(x,y)=P(x,y)u(y)/(Pu)(x)$ has stationary density $w$, which satisfies
the same floor, critical cap and exact matching conditions. Its eigenvalues,
ordered decreasingly with multiplicities, satisfy
\[
 |\lambda_j(K_u)-\lambda_j(P)|\leq\mathcal E_k
       \quad(1\leq j\leq2^L),\qquad
 |\operatorname{gap}(K_u)-(1-\rho)|\leq\mathcal E_k.
\]
The error decays exponentially in $k$, uniformly in $L,\pi,\tau$ for
fixed $\rho,b$. At $k=k_L$ in \eqref{admissible} it decays exponentially
in $\sqrt L$ for each fixed $C>0$.
\end{theorem}

\begin{proof}
\emph{Reconstruction and a derived stronger moment.}
Stochasticity, Jensen's inequality and the Walsh multipliers give
$\E w=1$, $w\geq b$, $\E w^s\leq2$ and the same zero coefficients
through degree $k$. Lemma~\ref{inversestationary} supplies the unique
positive solution of $uPu=w$, with $\E uPu=1$. Its dimension-dependent
stability estimate is not used here. Ordinary hypercontractivity
\cite{mos,odonnell} and $\tau^2\leq(s-1)/(\sigma-1)$ give
\begin{equation}\label{densityhighermoment}
 \norm w_\sigma\leq\norm\pi_s\leq B.
\end{equation}
This stronger norm bound is derived, not an additional premise. It does
not assert $\E w^\sigma\leq2$.

\emph{A finite inverse-moment crossing.}
Put $g=Pu$ and $Q=P_{\sqrt\rho}$. Then $\norm{Qu}_2=1$ and
$\norm g_R\leq1$ for $R=1+\rho^{-1}$, by ordinary
hypercontractivity applied to $g=Q(Qu)$. The floor $ug\geq b$ gives
$u^{-1}\leq g/b$, or $\norm u_{-R}\geq b$ in generalized-mean
notation. Reverse hypercontractivity, including its negative-exponent
range~\cite{mos}, applies since
$\rho^2=(1+R)/(1+q_0)$. Thus
$\norm{g^{-1}}_{q_0}\leq b^{-1}$.
All functions are strictly positive on a finite space.
The floor and rowwise Jensen also give
\[
 g\geq bP(1/g)\geq b/(Pg),\qquad g^{-1}\leq Pg/b.
\]
If $\norm{g^{-1}}_q\leq A_q$ and $q>\sigma/(\sigma-1)$,
H\"older with $p=\sigma q/(\sigma+q)>1$ gives
$\norm u_p=\norm{w/g}_p\leq BA_q$.
Since $Pg=P_{\rho^2}u$, ordinary hypercontractivity then yields
\begin{equation}\label{densitybootstrap}
 \norm{g^{-1}}_{F_\sigma(q)}\leq(B/b)A_q.
\end{equation}
This iteration uses the newly derived exponent $\sigma$, not the original
critical exponent's fixed-point formula.

For $q_0\leq q<Q_*$, direct algebra and the choice of $\sigma$ give
\begin{align*}
 F_\sigma(q)-q
 &=\frac{(1-\rho^4)(q-1)-q^2/(\sigma+q)}{\rho^4}\\
 &\geq\frac{(1-\rho^4)(q_0-1)}{2\rho^4}=:d>0.
\end{align*}
Here $q_0>2$, while $\sigma\geq4s>4$, so every input before the
first crossing satisfies $q>\sigma/(\sigma-1)$. If $q_0<Q_*$, absence
of a crossing would give $q_n\geq q_0+nd$ for all $n$, a contradiction.
In fact $n_*\leq\lceil(Q_*-q_0)/d\rceil$. If $q_0\geq Q_*$ no
step is needed. At the first crossing the last valid step has already
proved
\[
 \norm{g^{-1}}_{q_{n_*}}\leq A_*,\qquad q_{n_*}\geq8s.
\]
No assertion about further iterates is required.
For $U=u/g=w/g^2$, $N=\E U$ and $r=U/N$, H\"older gives
\begin{equation}\label{densityratiomoment}
 \norm U_{2s}\leq\norm w_{4s}\norm{g^{-1}}_{8s}^2
       \leq BA_*^2,\qquad \norm r_{2s}\leq C_*,
       \qquad \norm w_{2s}\leq B,
\end{equation}
where $N\geq b$ is justified independently in the next paragraph.

\emph{Actual normalization and matching rigidity.}
Let $a=\E u=\E g$ and $v=1-a^2=\norm{Qu-a}_2^2$.
The floor gives $a\geq b\E(1/g)\geq b/a$, so $a^2\geq b$.
Cauchy--Schwarz under $w\,d\mu$ gives
$N=\E_w(g^{-2})\geq(\E_w(g^{-1}))^2=a^2$.
For the other bound, the reversible conjugate
$S=M_{\sqrt U}PM_{\sqrt U}$ has Perron eigenvalue one and positive
spectrum. With $X=M_{\sqrt U}Q$, the matrices $S=XX^*$ and
$T=X^*X=QM_UQ$ are similar because $X$ is invertible. Thus
$\norm T_{2\to2}=1$ and testing at the unit constant gives $N\leq1$.
This proves $b\leq a^2\leq N\leq1$ without any perturbation estimate.

For completeness, the matching argument of Lemma~\ref{rigidity} can
be used directly with exponent $\sigma$. Since $\sigma'<2$, the
degree-$>k$ part of $g=Q(Qu)$ has $L^{\sigma'}$ norm at most
$\rho^{(k+1)/2}\sqrt v$. Exact matching gives
$\E w\Pi_{\leq k}g=a$, hence
$\E wg\leq a+e_k\sqrt v$.
Cauchy--Schwarz under $u\,d\mu$ gives
$1=(\E ug)^2\leq a\E ug^2=a\E wg$.
Consequently $v\leq ae_k\sqrt v$, including $v=0$, and
$a^2\geq(1+e_k^2)^{-1}$.
The actual identity $\E\sqrt{rw}=a/\sqrt N\geq a$ yields
\[
 \norm{\sqrt r-\sqrt w}_2^2
       \leq2(1-a)\leq e_k^2,\qquad \norm{r-w}_1\leq2e_k.
\]
Interpolation with \eqref{densityratiomoment} gives
\begin{equation}\label{densitycriticalcloseness}
 \norm{r-w}_s\leq(2e_k)^\gamma(C_*+B)^{1-\gamma},
 \qquad \gamma=\frac1{2s-1}.
\end{equation}

\emph{Closeness to the constant and the full spectrum.}
Factor $P_\tau=P_{\theta_{\rm out}}P_\beta P_{\theta_{\rm in}}$,
where $\beta=\tau/(\theta_{\rm in}\theta_{\rm out})\leq1/2$.
The input map contracts $L^s$ to $L^2$: use hypercontractivity if
$s<2$, and the identity with norm monotonicity if $s\geq2$.
The output contracts $L^2$ to $L^s$: use hypercontractivity if $s>2$,
and the identity with norm monotonicity if $s\leq2$.
The Fourier support of $\pi-1$ lies above degree $k$, so the middle
map contracts its $L^2$ norm by $\beta^{k+1}$. Therefore
\[
 \norm{w-1}_s\leq(B+1)2^{-(k+1)}.
\]
This also holds at $\tau=0$, where the mean-zero part vanishes.
Using $U-1=N(r-w)+N(w-1)+(N-1)$, $N\leq1$ and
$1-N\leq e_k^2/(1+e_k^2)\leq e_k^2$ gives
$\norm{U-1}_s\leq\mathcal E_k$.
The critical sandwich estimate of Lemma~\ref{criticalsandwich} gives
\[
 \norm{T-P}_{2\to2}=\norm{QM_{U-1}Q}_{2\to2}
              \leq\norm{U-1}_s\leq\mathcal E_k.
\]
Indeed, $2s'=1+\rho^{-1}$ and $Q:L^2\to L^{2s'}$ is a contraction,
so H\"older bounds $|\E(U-1)(Qf)(Q\ell)|$ by $\norm{U-1}_s$ for
arbitrary unit $L^2$ functions. This controls the full operator.
The self-adjoint min--max principle compares every ordered eigenvalue
of $T$ and $P$; the exact similarity above identifies the spectrum of
$T$ with that of $K_u$. No uniform condition number for $X$ is used.
Since $\lambda_2(P)=\rho$ for $L\geq1$, the gap bound follows.
All three terms of \eqref{densityspectrumerror} decay exponentially in
$k$ for fixed $\rho,b$. At $k=L$ the high-degree space is zero; at
$k=0$ the same argument uses mean one. These endpoints do not require
a separate assumption.
\end{proof}

\begin{proposition}[Fixed smoothing need not be a small density perturbation]\label{densityperturbationfailure}
Fix $C>0$ and $0\leq\tau<1$. There are energy-normalized positive
potentials $h_L$ satisfying all of \eqref{admissible} such that
\[
 \norm{\pi_L-P_\tau\pi_L}_s\longrightarrow\varepsilon>0.
\]
\end{proposition}
\begin{proof}
Take $\varepsilon=\min\{(1-b)/2,(B-1)/2\}>0$. For all sufficiently
large $L$, one has $k_L<L$. Set $\pi_L=1+\varepsilon\chi_{[L]}$ in
those dimensions. It has mean one, minimum at least $(1+b)/2>b$,
maximum at most $(1+B)/2<B$, and no nonconstant Fourier coefficients
below degree $L$. In particular $\E\pi_L^s<2$ and the required matching
holds. Lemma~\ref{inversestationary} realizes it as $h_LPh_L$ for a
positive energy-normalized potential. Use $h_L=1$ in the remaining
finitely many dimensions. Since $|\chi_{[L]}|=1$,
\[
 \norm{\pi_L-P_\tau\pi_L}_s
       =\varepsilon(1-\tau^L)\longrightarrow\varepsilon.
\]
\end{proof}
This refutes only the uniformly vanishing $L^s$ premise for a transfer
argument. It does not refute a spectral comparison: the parity family
itself belongs to an already controlled code class. Theorem~\ref{densityspectrumtheorem}
establishes the regularized spectrum on its specified interval, which
excludes $\tau=1$. Neither a comparison with the original spectrum nor
an extension to every $\tau<1$ is inferred.
